\chapter{Demonstration and Testing}

This chapter first demonstrates the completed platform through its core user workflows, then presents the testing and evaluation results that validate system correctness, performance, and usability.

% =========================================================
\section{System Demonstration}

This section walks through the platform's key features following the natural user journey, from onboarding to advanced risk analysis.

% ---------------------------------------------------------
\subsection{Registration}

New users register with an optional display name, email, and password. A confirm-password field with a real-time match indicator validates input. Upon successful registration, the user is redirected to the risk questionnaire to begin personalisation.

% ---------------------------------------------------------
\subsection{Risk Profiling}

The risk questionnaire comprises eight multiple-choice questions covering loss tolerance, time horizon, return expectations, reaction to downturns, income stability, liquidity needs, and portfolio goals. A progress bar tracks completion, and the ``Get Results'' button activates only once all questions are answered.

\begin{figure}[H]
\centering
\includegraphics[width=0.85\textwidth]{fig/demo_risk_quiz.png}
\caption{Risk profiling questionnaire with progress tracking}
\label{fig:demo_risk_quiz}
\end{figure}

The system computes a weighted score and classifies the user as \textit{Conservative}, \textit{Balanced}, or \textit{Aggressive}. The result screen shows the profile badge, total score, an intent snapshot (timeline, goal, drawdown comfort, cash priority), and a per-question breakdown. The user can then view their top ETF picks or edit their answers.

\begin{figure}[H]
\centering
\includegraphics[width=0.85\textwidth]{fig/demo_risk_result.png}
\caption{Risk quiz result with profile classification and intent snapshot}
\label{fig:demo_risk_result}
\end{figure}

% ---------------------------------------------------------
\subsection{Personalised ETF Recommendations}

The recommendations page displays five ETF cards ranked by the multi-factor score, pre-filtered to the user's risk profile. Each card shows the ticker, name, category, key metrics (YTD return, beta, expense ratio, dividend yield), a score out of 100, and reason labels explaining the match (e.g., ``Low volatility,'' ``Strong dividend yield''). Users can switch profiles to compare recommendations, click a card to view its detail page, or navigate to the full ETF discovery screen.

\begin{figure}[H]
\centering
\includegraphics[width=0.85\textwidth]{fig/demo_recommendations.png}
\caption{Personalised ETF recommendations with scores and explanatory labels}
\label{fig:demo_recommendations}
\end{figure}

% ---------------------------------------------------------
\subsection{ETF Discovery}

The ETF Discovery page provides access to all 2,272~ETFs via a search bar and filters for category, asset class, ETF type, region, YTD direction, beta profile, and expense ratio. Results appear in a sortable, paginated table. A collapsible ``ETF Quick Guide'' explains key concepts for beginner users.

\begin{figure}[H]
\centering
\includegraphics[width=0.85\textwidth]{fig/demo_etf_discovery.png}
\caption{ETF Discovery page with search, filters, and sortable table}
\label{fig:demo_etf_discovery}
\end{figure}

% ---------------------------------------------------------
\subsection{ETF Detail View}

Selecting an ETF opens its detail page, which includes a quick-stats bar (YTD return, expense ratio, AUM, beta, dividend yield, 52-week range), an interactive price chart (1M--MAX), computed risk metrics (volatility, beta, Sharpe ratio, max drawdown), holdings breakdown, fund metadata, and a ticker-level news feed with sentiment scores.

\begin{figure}[H]
\centering
\includegraphics[width=0.85\textwidth]{fig/demo_etf_detail.png}
\caption{ETF detail page with price chart, risk metrics, and news feed}
\label{fig:demo_etf_detail}
\end{figure}

% ---------------------------------------------------------
\subsection{AI Chatbot Interaction}

Each detail page includes an AI chatbot accessible via a floating button. The chatbot is grounded with the ETF's live metadata and metrics, enabling factual responses to questions such as ``Is this ETF suitable for a conservative investor?'' Every reply includes a financial disclaimer. Figure~\ref{fig:demo_chatbot} shows a sample conversation.

\begin{figure}[H]
\centering
\includegraphics[width=0.85\textwidth]{fig/demo_chatbot.png}
\caption{AI chatbot conversation grounded with live ETF data}
\label{fig:demo_chatbot}
\end{figure}

% ---------------------------------------------------------
\subsection{Wallet Management}

Before adding holdings, users set up wallets to define their investment goals (e.g., ``Retirement,'' ``Education,'' ``Growth''). Each wallet has configurable profile settings---risk profile, investment horizon, objective, liquidity preference, and maximum drawdown---allowing distinct strategies within one account. Wallets are displayed as cards showing purpose, invested value, market value, and P/L.

\begin{figure}[H]
\centering
\includegraphics[width=0.85\textwidth]{fig/demo_wallets.png}
\caption{Wallet management with purpose-based organisation and per-wallet profiles}
\label{fig:demo_wallets}
\end{figure}

% ---------------------------------------------------------
\subsection{Portfolio Management}

With wallets in place, users add ETF holdings via the ``Add to Portfolio'' button on any detail page or wallet page. The modal collects the ticker, quantity, purchase price, date, and target wallet. The portfolio overview aggregates all holdings and displays total market value, invested capital, unrealised P/L, a growth chart, an allocation pie chart, and a holdings table with inline editing. A wallet filter scopes the view to a specific wallet, and a news section surfaces articles for held tickers.

\begin{figure}[H]
\centering
\includegraphics[width=0.85\textwidth]{fig/demo_portfolio.png}
\caption{Portfolio overview with summary, growth chart, allocation, and holdings}
\label{fig:demo_portfolio}
\end{figure}

% ---------------------------------------------------------
\subsection{Scenario Analysis}

The Scenario Analysis page lets users stress-test holdings at the portfolio, wallet, or single-ETF level. Available scenarios include a COVID-19 crash replay (2020-02-19 to 2020-03-23) and uniform drawdowns ($-10\%$, $-20\%$, $-30\%$). Results show baseline versus scenario values, percentage change, and per-holding impact with worst/best performer insights.

\begin{figure}[H]
\centering
\includegraphics[width=0.85\textwidth]{fig/demo_scenario.png}
\caption{COVID-19 scenario analysis with per-holding impact breakdown}
\label{fig:demo_scenario}
\end{figure}

% ---------------------------------------------------------
\subsection{VaR/CVaR Tail Risk Analysis}

A second tab provides VaR and CVaR at configurable confidence levels (90\%, 95\%, 99\%) and time horizons (1M, 3M, 1Y, 2Y) using both historical and parametric methods. Results include dollar-amount risk estimates and portfolio statistics (mean daily return, volatility, worst/best day).

\begin{figure}[H]
\centering
\includegraphics[width=0.85\textwidth]{fig/demo_var.png}
\caption{VaR/CVaR tail risk analysis with historical and parametric estimates}
\label{fig:demo_var}
\end{figure}




% =========================================================
\section{Testing and Evaluation}

Having demonstrated the platform's features, this section presents the testing methodology and results that validate system correctness, performance, and usability. Backend tests were executed using FastAPI TestClient and Swagger UI against a development instance connected to the Supabase PostgreSQL database.

\subsection{API Testing}

The backend exposes 49 endpoints in total (47 router endpoints across 12 modules and 2 system endpoints in \texttt{main.py}: \texttt{/} and \texttt{/health}). A full API validation cycle was executed across all endpoints, with 124 consolidated test cases covering success paths, input validation, authentication enforcement, and error handling.

\begin{table}[H]
\centering
\begin{tabular}{lccc}
\toprule
\textbf{Module} & \textbf{Endpoints} & \textbf{Test Cases} & \textbf{Pass Rate} \\
\midrule
Authentication & 4 & 14 & 100\% \\
ETF Data (etfs + prices) & 7 & 18 & 100\% \\
Metrics Engine & 3 & 10 & 100\% \\
News & 6 & 14 & 100\% \\
Macro Indicators & 4 & 10 & 100\% \\
Risk Profiling (quiz) & 3 & 10 & 100\% \\
Recommendations & 2 & 8 & 100\% \\
Portfolio & 5 & 14 & 100\% \\
Wallet Management & 8 & 16 & 100\% \\
Scenario Analysis & 4 & 6 & 100\% \\
Chatbot & 1 & 2 & 100\% \\
System (root / health) & 2 & 2 & 100\% \\
\midrule
\textbf{Total} & \textbf{49} & \textbf{124} & \textbf{100\%} \\
\bottomrule
\end{tabular}
\caption{API Endpoint Test Coverage}
\label{tab:api_test_summary}
\end{table}

Test cases included authentication checks (duplicate signup rejected with 400, missing token returns 401, cross-user wallet access returns 403), data retrieval validation (invalid ticker returns 404, paginated ETF list returns correct page size), input validation via Pydantic schemas (type mismatches and missing fields return 422), and computation verification (VaR/CVaR returns both historical and parametric values, scenario analysis returns per-holding impact).

% =========================================================
\subsection{Performance Evaluation}

System performance was benchmarked against the non-functional requirements defined in Chapter~3.

\begin{table}[H]
\centering
\begin{tabular}{llll}
\toprule
\textbf{Operation} & \textbf{Metric} & \textbf{Target} & \textbf{Result} \\
\midrule
ETF List Load (paginated) & Response Time & $<$500ms & 120ms \\
ETF Detail with Metrics & Response Time & $<$500ms & 280ms \\
Scenario Analysis (COVID) & Calculation Time & $<$3000ms & 1200ms \\
VaR/CVaR (252-day, 95\%) & Calculation Time & $<$3000ms & 950ms \\
AI Chatbot & Response Generation & $<$5000ms & 2500ms \\
Recommendation Generation & Response Time & $<$1000ms & 350ms \\
Lighthouse Score & Performance & $>$80 & 88 \\
\bottomrule
\end{tabular}
\caption{Performance Benchmarks}
\label{tab:performance_benchmarks}
\end{table}

All operations met or exceeded their target thresholds. The chatbot response time is dominated by the external OpenAI API call latency rather than internal processing.

% =========================================================
\subsection{Frontend Testing}

Google Lighthouse was run on the deployed production frontend (Vercel) using the ETF Discovery page as a representative data-heavy page.

\begin{table}[H]
\centering
\begin{tabular}{llll}
\toprule
\textbf{Category} & \textbf{Target} & \textbf{Result} & \textbf{Rating} \\
\midrule
Performance & $>$80 & 88 & Good \\
Accessibility & $>$80 & 91 & Good \\
Best Practices & $>$80 & 95 & Good \\
SEO & $>$80 & 92 & Good \\
First Contentful Paint & $<$1.5s & 0.8s & Good \\
Largest Contentful Paint & $<$2.5s & 1.6s & Good \\
\bottomrule
\end{tabular}
\caption{Lighthouse Audit Results}
\label{tab:lighthouse}
\end{table}

Responsive design was verified across mobile (375$\times$812), tablet (768$\times$1024), and desktop (1440$\times$900) viewports using Chrome DevTools. Core user flows were also validated on Chrome, Firefox, Safari, and Edge without functional issues.

% =========================================================
\subsection{User Acceptance Testing (UAT)}

A UAT session was conducted with 8 retail investors (age 18-22; 5 beginners, 3 intermediate) to evaluate usability. Participants completed seven tasks without assistance:

\begin{table}[H]
\centering
\begin{tabular}{p{0.8cm} p{6.5cm} c c}
\toprule
\textbf{Task} & \textbf{Description} & \textbf{Completion} & \textbf{Avg.\ Time} \\
\midrule
T1 & Register and complete the risk questionnaire & 100\% & 2.5 min \\
T2 & View personalised ETF recommendations & 100\% & 1.0 min \\
T3 & Search for a specific ETF and view its details & 100\% & 1.5 min \\
T4 & Add two ETF holdings to the portfolio & 100\% & 2.0 min \\
T5 & Create a wallet and assign a holding & 100\% & 2.5 min \\
T6 & Run a scenario analysis (COVID crash) on portfolio & 87.5\% & 3.0 min \\
T7 & Ask the AI chatbot a question about an ETF & 100\% & 1.0 min \\
\bottomrule
\end{tabular}
\caption{UAT Task Completion Results}
\label{tab:uat_tasks}
\end{table}

Task T6 had one incomplete attempt where the participant initially could not locate the scenario analysis page; the task was completed after brief orientation, prompting a navigation label improvement. Users rated Ease of Navigation at 4.5/5, Usefulness of Recommendations at 4.3/5, Clarity of Risk Metrics at 4.1/5, Helpfulness of AI Chatbot at 4.4/5, and Overall Satisfaction at 4.4/5, validating the platform's core value proposition.
